\section{Related Work}
\label{sec:related}

\paragraph{Cybersecurity Benchmarks for AI Agents}
We now compare \sysname with recent cybersecurity benchmarks, as detailed in \cref{tab:comp}.
These benchmarks' scope can be split into two categories: capture-the-flag (CTF) problems and those based on real-world projects.
Earlier benchmarks like NYU CTF Bench~\citep{shao2024nyu} and \cybench~\citep{cybench} rely exclusively on CTF problems.
Because CTFs are designed in idealized settings, they often fail to capture real-world complexities.
Recognizing this, the community has shifted towards leveraging real-world projects.
This includes AutoAdvExBench~\citep{carlini2025autoadvexbench}, CVE-Bench~\citep{zhu2025cve}, BountyBench~\citep{zhang2025bountybench}, SEC-Bench~\citep{lee2025sec}, and our own \sysname.

\sysname stands out in both scale and diversity.
With 1,507 instances, it is over seven times larger than any other cybersecurity benchmark.
Furthermore, these instances are derived from 188 software projects from diverse application domains, as listed in \cref{tab:all_projects}.
This ensures that \sysname effectively measures progress by capturing a wide range of difficulties, as demonstrated by the gradually improved performance of frontier models in our evaluation (\cref{sec:eval}).

Another key differentiator for our work is its in-depth analysis on agents' ability to discover new, zero-day vulnerabilities (\cref{sec:discovery}).
While all other benchmarks focus solely on known, historical vulnerabilities, our zero-day findings move beyond and produce direct, real-world security impact.

\begin{table}[h]
\centering
\vspace{-2mm}
\caption{Comparing \sysname with existing cybersecurity benchmarks for AI agents.}
\label{tab:comp}
\vspace{-2mm}
\resizebox{\textwidth}{!}{
\begin{tabular}{@{}lrrrr@{}}
\toprule
\textbf{Benchmark} & \textbf{Scope} & \textbf{\# Instances} & \textbf{\# Projects} & \textbf{Zero-days} \\
\midrule
NYU CTF Bench~\citep{shao2024nyu} & CTF & 200 & - & \xmark \\
\cybench~\citep{cybench} & CTF & 40 & - & \xmark \\
AutoAdvExBench~\citep{carlini2025autoadvexbench} & CTF+Real-world & 75 & 41 & \xmark \\
CVE-Bench~\citep{zhu2025cve} & Real-world & 40 & 26 & \xmark \\
BountyBench~\citep{zhang2025bountybench} & Real-world & 40 & 31 & \xmark \\
SEC-bench~\citep{lee2025sec} & Real-world & 200 & 29 & \xmark \\
\midrule
\sysname{} (Our work) & Real-world & 1,507 & 188 & \cmark \\
\bottomrule
\end{tabular}
}
\vspace{-2mm}
\end{table}


\paragraph{Coding Benchmarks for AI Agents}
Existing coding benchmarks such as SWE-bench~\citep{jimenez2024swebench} and SWT-bench~\citep{mundler2024swtbench} evaluate AI agents' ability to handle software engineering tasks.
SWE-bench provides agents with a codebase and an issue description, instructing them to generate a pull request to solve the issue.
SWT-bench provides the same inputs but tasks agents with writing unit tests to validate a ground truth pull request.
These benchmarks have sparked the development of various coding agents, such as OpenHands~\citep{openhands} and Codex~\citep{openai_codex}, as well as specialized backbone models like SWE-Gym~\citep{swegym} and R2E-Gym~\citep{r2egym}, which are fine-tuned to achieve high performance on SWE-bench.

While \sysname can be seen as a coding benchmark, it focuses specifically on security, in contrast to the functionality-focused nature of SWE-bench and SWT-bench.
SWE-bench and SWT-bench often involve making localized code changes, whereas \sysname requires more comprehensive, repository-wide reasoning.
To succeed on \sysname, an agent must craft a proof of concept input that accurately navigates from the program's entry point to the vulnerability, demanding a deep understanding of the entire codebase.
Due to these differences, general-purpose software agents and LLMs specially fine-tuned for software engineering tasks struggle on \sysname, as evidenced by our evaluation results in \cref{sec:eval}.
This highlights \sysname's complementary value to existing coding benchmarks such as SWE-bench and its importance for a more complete agent evaluation.
